\section{Sesión de laboratorio: Lab de detección de rostros}

Amini presentó en clase el laboratorio complementario de esta clase (Lab 2): la tarea de detección de rostros. Este laboratorio permite a los estudiantes construir con sus propias manos una CNN para detectar rostros en imágenes, una excelente oportunidad para llevar a la práctica la teoría de esta clase.

Los pasos principales del laboratorio incluyen:
\begin{enumerate}
    \item Construir la arquitectura de la red neuronal convolucional (incluyendo capas convolucionales, activación ReLU, capas de pooling y una cabeza de clasificación totalmente conectada)
    \item Entrenar el modelo con un conjunto de datos de imágenes de rostros
    \item Evaluar el desempeño de detección del modelo en imágenes nuevas
    \item Analizar la visualización de las características aprendidas por el modelo
\end{enumerate}

\begin{knowledgebox}{El significado histórico de la detección de rostros}
La detección de rostros es una de las primeras tareas de visión por computadora en obtener una aplicación comercial a gran escala. Del detector de rostros Viola-Jones (2001, basado en características Haar diseñadas manualmente y en un clasificador AdaBoost) hasta los métodos modernos de aprendizaje profundo, la precisión y la velocidad de la detección de rostros han dado un salto cualitativo. Hoy en día, la tecnología de detección de rostros se ha integrado en aplicaciones cotidianas como el desbloqueo de teléfonos, los filtros de redes sociales y la vigilancia de seguridad.
\end{knowledgebox}

\begin{warningbox}{Consideraciones éticas del reconocimiento facial}
Aunque las tecnologías de detección y reconocimiento facial son muy poderosas, también han suscitado graves problemas de privacidad y ética. Estos incluyen: sesgos raciales y de género (el modelo se desempeña peor en ciertos grupos), abuso de la vigilancia y falsificaciones mediante deepfake, entre otros. Al desarrollar y desplegar este tipo de tecnología, es indispensable considerar seriamente su impacto social.
\end{warningbox}

\newpage
